\documentclass[11pt, letterpaper]{article}

\usepackage[top=0.85in, bottom=0.85in, left=0.85in, right=0.85in]{geometry}
\usepackage{amsmath, amssymb, amsthm, mathtools}
\usepackage{longtable}
\usepackage{booktabs}
\usepackage{array}
\usepackage{tabularx}
\usepackage{xcolor}
\usepackage{fancyhdr}
\usepackage{titlesec}
\usepackage{setspace}
\usepackage[T1]{fontenc}
\usepackage{enumitem}
\usepackage{tocloft}
\usepackage{tcolorbox}
\tcbuselibrary{breakable, skins}

% --- Color Definitions ---
\definecolor{RSblue}{RGB}{25, 60, 120}
\definecolor{RSgray}{RGB}{90, 90, 90}
\definecolor{RSgreen}{RGB}{20, 110, 60}
\definecolor{RSred}{RGB}{160, 30, 30}
\definecolor{RSamber}{RGB}{180, 110, 0}
\definecolor{RSpurple}{RGB}{90, 30, 120}

% --- Box Environments ---
\newtcolorbox{rsbox}{
  enhanced, breakable,
  colback=RSblue!4, colframe=RSblue!60,
  title={\textbf{RS Ontological Statement}},
  fonttitle=\small\bfseries\color{white},
  attach boxed title to top left={yshift=-2mm, xshift=4mm},
  boxed title style={colback=RSblue},
  left=4pt, right=4pt, top=6pt, bottom=4pt
}

\newtcolorbox{verifybox}{
  enhanced, breakable,
  colback=RSgreen!4, colframe=RSgreen!50,
  title={\textbf{External Verification (Continuous Lens)}},
  fonttitle=\small\bfseries\color{white},
  attach boxed title to top left={yshift=-2mm, xshift=4mm},
  boxed title style={colback=RSgreen!70!black},
  left=4pt, right=4pt, top=6pt, bottom=4pt
}

\newtcolorbox{observebox}{
  enhanced, breakable,
  colback=RSamber!5, colframe=RSamber!60,
  title={\textbf{Structural Observation}},
  fonttitle=\small\bfseries\color{white},
  attach boxed title to top left={yshift=-2mm, xshift=4mm},
  boxed title style={colback=RSamber!80!black},
  left=4pt, right=4pt, top=6pt, bottom=4pt
}

% --- Page Style ---
\pagestyle{fancy}
\fancyhf{}
\fancyhead[L]{\small\textcolor{RSgray}{RigbySpace Discrete Dynamics}}
\fancyhead[R]{\small\textcolor{RSgray}{Rational Euler Dynamics -- Page \thepage}}
\fancyfoot[L]{\small\textcolor{RSgray}{D. Veneziano}}
\fancyfoot[R]{\small\textcolor{RSgray}{Discrete Unification Framework}}
\renewcommand{\headrulewidth}{0.4pt}

% --- Section Formatting ---
\titleformat{\section}
  {\large\bfseries\color{RSblue}}{\thesection}{1em}{}[\titlerule]
\titleformat{\subsection}
  {\normalsize\bfseries\color{RSblue}}{\thesubsection}{1em}{}
\titleformat{\subsubsection}
  {\normalsize\bfseries\color{RSgray}}{\thesubsubsection}{1em}{}

\onehalfspacing
\setlength{\parskip}{4pt}
\setlength{\parindent}{0pt}

% --- Macros ---
\newcommand{\Z}{\mathbb{Z}}
\newcommand{\SL}{S_{\mathrm{L}}}
\newcommand{\koppa}{\mathbin{\text{\rm\k{o}}}}
\newcommand{\succOp}{\operatorname{succ}}
\newcommand{\PG}{\mathrm{PG}}
\newcommand{\Dcross}{\Delta_{\mathrm{cross}}}
\newcommand{\Tvac}{T_{\mathrm{vac}}}
\newcommand{\Tbary}{T_{\mathrm{bary}}}
\newcommand{\Dg}{D_{\mathrm{gauge}}}
\newcommand{\Sint}{S_{\mathrm{int}}}
\newcommand{\Sigimb}{\Sigma_{\mathrm{imb}}}
\newcommand{\Gmg}{G_{\mathrm{mg}}}
\newcommand{\dslip}{\delta_{\mathrm{slip}}}
\newcommand{\tauX}{\tau_{10}}
\newcommand{\aSC}{\alpha^{-1}_{\mathrm{int}}}
\newcommand{\med}{\oplus}
\newcommand{\aop}{\eta}
\newcommand{\lop}{\lambda}
\newcommand{\psiop}{\psi}
\newcommand{\rank}{\rho}
\newcommand{\Sigs}{\Sigma_{\mathrm{s}}}
\newcommand{\Xis}{\Xi_{\mathrm{s}}}
\newcommand{\Omegas}{\Omega_{\mathrm{s}}}
\newcommand{\Pithree}{\Pi_3}
\newcommand{\nG}{n_G}
\newcommand{\Ksat}{K_{\mathrm{sat}}}
\newcommand{\RD}{R_D}
\newcommand{\bplus}{\boxplus}

\begin{document}

\begin{titlepage}
  \centering
  \vspace*{2cm}
  {\Huge\bfseries\color{RSblue} RigbySpace Discrete Dynamics\par}
  \vspace{0.5cm}
  {\LARGE Rational Euler Dynamics\par}
  \vspace{0.5cm}
  {\large Discrete Emulation of Complex Rotation and Wave Mechanics\par}
  \vspace{2cm}
  \rule{0.6\textwidth}{0.4pt}\par
  \vspace{1cm}
  {\large D.\ Veneziano\par}
  \vspace{0.8cm}
  {\normalsize RigbySpace Discrete Dynamics\par}
  \vspace{1.5cm}
  \rule{0.6\textwidth}{0.4pt}\par
  \vspace{0.5cm}
  {\small $\Tvac = 11 \quad N = 3$ \par}
  \vfill
  {\small\color{RSgray}
  RigbySpace is a formally consistent candidate theory internal to its own
  axioms. No empirical equivalence to external physical theory is asserted
  unless an explicit mapping is provided and derived under that mapping.\par}
\end{titlepage}

\tableofcontents
\newpage

%=====================================================================
\section{Introduction: The Illusion of the Complex Plane}
%=====================================================================

In continuous mathematics and quantum mechanics, rotations, oscillations, and wave functions are elegantly described using complex numbers and Euler's formula ($e^{i\theta} = \cos\theta + i\sin\theta$). The imaginary unit $i$ acts as a continuous $90^\circ$ rotation operator in an abstract complex plane.

RigbySpace (RS) strictly forbids complex numbers, transcendental functions, and continuous geometry from its ontological chain. Therefore, the wave-like behavior of particles and the rotational dynamics of spinors must emerge entirely from discrete integer operations on unreduced rational pairs.

This document provides the complete, granular formalization of Rational Euler Dynamics. It demonstrates how a coupled system of two integer streams (a "real" and "imaginary" register) interacts through Surge operators ($\aop, \lop$) and the Transformative Reciprocal ($\psiop$) to emulate rotation, damped harmonic oscillation, and waveform synthesis. It proves that the imaginary unit $i$ is simply the macroscopic, continuous mask of the $\psiop$ operator exchanging unreduced causal history between two states.

%=====================================================================
\section{The Coupled State and the Discrete $90^\circ$ Phase Shift}
%=====================================================================

To emulate a 2D complex plane without invoking $\mathbb{C}$, we utilize a coupled pair of Explicit Rational Pairs (ERPs).

\subsection{The Rational State Pair}

\begin{rsbox}
\begin{definition}[Rational State Pair]\label{def:state_pair}
A rational state at causal tick $k$ is an ordered pair of ERPs, $S_k = (A_k, B_k)$, where:
\[ A_k = (n_a, d_a) \in \Z_{>0}^2 \]
\[ B_k = (n_b, d_b) \in \Z_{>0}^2 \]

*   The register $A_k$ is designated the \textbf{Real Component}.
*   The register $B_k$ is designated the \textbf{Imaginary Component}.

No cancellation of factors (GCD reduction) occurs. The denominators $d_a$ and $d_b$ record the entire structural history of operations applied to their respective registers.
\end{definition}
\end{rsbox}

\subsection{The Transformative Reciprocal as the Imaginary Unit}

In continuous math, multiplying by $i$ rotates a vector by $90^\circ$. Multiplying by $i^2 = -1$ inverts it ($180^\circ$). Multiplying by $i^4 = 1$ returns it to its original state ($360^\circ$). In RigbySpace, this exact algebraic structure is generated by the Transformative Reciprocal.

\begin{theorem}[$\psiop$ as the Discrete Phase Shift]\label{thm:psi_rotation}
\begin{rsbox}
The Transformative Reciprocal $\psiop$ acts on the coupled state $S = (A, B)$ by:
\[ \psiop(A, B) = \psiop\left(\frac{n_a}{d_a}, \frac{n_b}{d_b}\right) = \left(\frac{d_b}{n_a}, \frac{d_a}{n_b}\right) \]

This operator swaps the structural histories (denominators) between the two registers, while converting the original denominators into new numerators. This violent exchange of inertia forces a discrete quarter-turn ($90^\circ$ phase shift) in the geometry of the state.

\textbf{The Four Iterates (The Spinor Cycle):}
1. \textbf{Quarter Turn ($90^\circ$):} $\psiop^1(A, B) = (d_b/n_a,\; d_a/n_b)$
2. \textbf{Half Turn ($180^\circ$):} $\psiop^2(A, B) = (n_b/d_b,\; n_a/d_a) = (B, A)$
   *(This exchanges the pair and inverts the cross-determinant sign, executing Charge Conjugation).*
3. \textbf{Three-Quarter Turn ($270^\circ$):} $\psiop^3(A, B) = (d_a/n_b,\; d_b/n_a)$
4. \textbf{Full Rotation ($360^\circ$):} $\psiop^4(A, B) = (n_a/d_a,\; n_b/d_b) = (A, B)$

The period of $\psiop$ is exactly 4. It is the discrete, integer-native equivalent of the imaginary unit $i$.
\end{rsbox}
\end{theorem}

%=====================================================================
\section{The Rational Euler Cycle}
%=====================================================================

We now assemble the Surge operators and the Phase Shift operator into a repeating cycle that emulates a rotating, oscillating vector.

\subsection{The Surge Operators}

\begin{rsbox}
\textbf{1. Aggregate Surge ($\aop$):}
\[ \aop(n, d) = (n + d, n) \]
Increases the numerator by the denominator, and replaces the denominator with the previous numerator. This forces the denominator to grow multiplicatively, injecting structural history (inertia) into the state.

\textbf{2. Linear Surge ($\lop$):}
\[ \lop(n, d) = (n + d, d) \]
Adds the denominator to the numerator while leaving the denominator unchanged. This generates a linear ramp of the kinematic magnitude.
\end{rsbox}

\subsection{The Four-Step Euler Sequence}

\begin{theorem}[The Rational Euler Cycle]\label{thm:euler_cycle}
\begin{rsbox}
Given an initial coupled state $S_0 = (A_0, B_0)$, one Rational Euler Cycle consists of the following four-step sequence:
\[ \mathcal{C}_{\text{Euler}} = [\aop_A, \psiop, \aop_B, \psiop] \]

\textbf{Step 1 (Real Surge):} $A \leftarrow \aop(A)$.
The Real register grows its structural history.
\textbf{Step 2 (Quarter Turn):} $(A, B) \leftarrow \psiop(A, B)$.
The massive new history of the Real register is injected into the Imaginary register.
\textbf{Step 3 (Imaginary Surge):} $B \leftarrow \aop(B)$.
The Imaginary register grows its structural history.
\textbf{Step 4 (Quarter Turn):} $(A, B) \leftarrow \psiop(A, B)$.
The history is swapped back, completing the cycle.
\end{rsbox}
\end{theorem}

%=====================================================================
\section{The Hidden Third Element and Damped Harmonic Oscillation}
%=====================================================================

The Rational Euler Cycle does not produce a static, perfect circle. It produces a damped harmonic oscillation. This damping is not an energy loss to an external environment; it is the deterministic consequence of the No-GCD axiom and the Thermodynamic Transfer Law.

\begin{theorem}[Denominator Growth and Monotonic Decay]\label{thm:damping}
\begin{rsbox}
Let $(n,d) \in \Z_{>0}^2$ with $n > d$. Applying the Aggregate Surge $\aop$ yields $(n+d, n)$. The new denominator $n$ is strictly greater than the original denominator $d$.

Once a register satisfies $n > d$, successive applications of $\aop$ maintain strict, multiplicative growth of its denominator. Because denominators are exchanged via $\psiop$ but never reduced, the structural history of the coupled system grows monotonically.

\textbf{The Damping Mechanism:}
Tracking one cycle shows that the Real register $A$ evolves from $(n_a, d_a)$ to $(n_a, n_a + n_b)$.
Because the Surge operations cause the denominators to grow massively, there exists a cycle index $k$ such that the denominator of the Real register permanently exceeds its numerator.
Beyond this point, each subsequent cycle reduces the rational value of the Real register ($v_A = n_a / d_a$), driving it monotonically toward zero.

Meanwhile, the Imaginary register $v_B$ oscillates between values determined by the alternating surge and twist.
\end{rsbox}
\end{theorem}

\begin{observebox}
\textbf{Thermodynamic Ground State Stabilization:}
The damping of the Rational Euler Cycle is the exact mechanical realization of the Thermodynamic Transfer Law derived in the Shell Dynamics sector:
\[ \chi = \varkappa + \varepsilon \]
(Total Inertial Load = Rest Inertia + Kinetic Excess).

The $\aop$ operator generates Kinetic Excess ($\varepsilon$) in the numerator. The $\psiop$ operator forces that excess into the denominator, converting it into Rest Inertia ($\varkappa$). The rational value of the state decays because the kinetic energy is being deterministically absorbed by the growing structural rank of the unreduced integers. The system naturally seeks the Ground State Equilibrium ($\varepsilon = 0$).
\end{observebox}

%=====================================================================
\section{Programmable Rational Waveform Synthesis}
%=====================================================================

Beyond the canonical Euler cycle, the TRTS engine can engineer distinct discrete waveforms by altering the sequence of Surge and Twist operations. This proves that all standard wave mechanics are macroscopic artifacts of specific integer operator loops.

\begin{theorem}[Waveform Programs]\label{thm:waveforms}
\begin{rsbox}
A \textbf{Waveform Program} is a finite ordered list of elementary operations $\{\aop_A, \aop_B, \lop_A, \lop_B, \psiop\}$. Executing the program sequentially constitutes one patch. Repeating the patch yields a discrete signal extracted from the Real register $v_A$.

\textbf{1. Damped Sine Wave (The Euler Cycle):}
*   Program: $[\aop_A, \psiop, \aop_B, \psiop]$
*   Result: The Real register displays a decaying envelope akin to a harmonic oscillator, spiraling toward a limit cycle in the $(v_A, v_B)$ plane.

\textbf{2. Square Wave (Charge and Hold):}
*   Program: $[\aop_A \times m, \psiop, \aop_B \times m, \psiop]$ (for integer $m > 0$).
*   Result: $\aop_A$ rapidly charges the Real register to a saturation limit (HIGH state). The $\psiop$ operator swaps the massive denominator to the numerator, instantly dropping the value (LOW state). Holding one register fixed while charging the other produces a discrete square wave with a duty cycle determined by $m$.

\textbf{3. Sawtooth Wave (Linear Ramp and Reset):}
*   Program: $[\lop_A \times m, \psiop]$
*   Result: The Linear Surge $\lop_A$ ramps the magnitude of the Real register linearly at a constant capacity. The $\psiop$ operator then inverts the state, providing a sharp, discrete reset. This is the exact integer mechanism of gauge boson propagation.

\textbf{4. Triangle Wave (Alternating Ramps):}
*   Program: $[\lop_A \times m, \psiop, \lop_B \times m, \psiop]$
*   Result: Alternates $m$ linear surges on A and $m$ on B between twists. Produces a wave that rises linearly, drops sharply, holds while the other register ramps, then rises again.
\end{rsbox}
\end{theorem}

%=====================================================================
\section{Conclusion}
%=====================================================================

Rational Euler Dynamics demonstrates that rotation, oscillation, and wave mechanics do not require the continuum, complex numbers, or trigonometric functions. 

By coupling two unreduced rational streams and subjecting them to a strict schedule of Aggregate Surges and Transformative Reciprocals, the RigbySpace engine natively generates the exact behaviors ordinarily attributed to complex analysis. The imaginary unit $i$ is revealed to be the macroscopic mask of the $\psiop$ operator exchanging unreduced causal history between registers. The damping of the harmonic oscillator is revealed to be the deterministic absorption of kinetic excess into structural inertia. 

The universe is not a smooth, continuous wave; it is a discrete, programmable integer synthesizer.

\end{document}
